% \section{Related Work}
\subsection{Prior approaches for code localization}
Several prior methods have developed methods targeting code localization for downstream issue resolution in repositories. \citet{jimenez2024swebenchlanguagemodelsresolve} identify relevant source code files using BM25 retrieval~\citep{article} treating the issue description as the query and the Python source code files as documents. SWE-Fixer~\citep{xie2025swefixertrainingopensourcellms} use a two-step coarse-to-fine localization approach where they first retrieve relevant source code files using BM25 retrieval and then prompt a fine-tuned 7B model with the skeleton of the retrieved files to predict the names of relevant files for issue resolution. On the other hand, Agentless~\citep{xia2024agentless} uses a complex, multi-step localization pipeline wherein it first retrieves suspicious code files from the repository using both embedding-based retrieval and by prompting an LLM with the high-level repository structure, followed by prompting an LLM to predict relevant functions and classes given the skeletons of the suspicious files. 

Many recent approaches have increasingly shifted towards developing models and systems for code localization using agentic scaffolds. LocAgent~\citep{chen-etal-2025-locagent} and RepoGraph~\citep{ouyang2025repographenhancingaisoftware} utilize a graph-based indexing approach wherein all dependencies in the code repository (for e.g. import, invoke, inherit, etc.) are captured in a code graph and also develop specialized tools allowing the agent to search and traverse the graph. Furthermore, \citet{chen-etal-2025-locagent} train 7B model for their scaffold using rejection sampling fine-tuning on successful trajectories sampled from Claude-3.5-Sonnet. On the other hand, CoSIL~\citep{liu2025cosil} leverages module-call and function-call graphs 

Recent LLM-based systems increasingly treat localization as an agentic process. Fixed-pipeline systems such as Agentless~\citep{xia2024agentless} separate localization from editing. Graph- and structure-aware methods, including LocAgent~\citep{chen-etal-2025-locagent}, CoSIL~\citep{liu2025cosil}, and RepoGraph~\citep{ouyang2025repographenhancingaisoftware}, introduce explicit repository structure to guide search. Other work explores tool-augmented or agent-based localization strategies, including priority scheduling and context pruning (e.g., OrcaLoca~\citep{yu2025orcalocallmagentframework}), as well as end-to-end RL training for repository-level agents such as RepoNavigator~\citep{zhang2025one}.

\subsection{Training methods for code agents}
Training code agents raises challenges in data efficiency, stability, and generalization. Many systems rely on supervised fine-tuning and/or distillation from strong teachers (e.g., LocAgent~\citep{chen-etal-2025-locagent} and SWE-Fixer~\citep{xie2025swe}), whereas reinforcement learning enables direct optimization from task rewards. Recent RL methods such as GRPO~\citep{shao2024deepseekmathpushinglimitsmathematical} and GSPO~\citep{zheng2025groupsequencepolicyoptimization} provide scalable policy optimization for language models. SWE-Gym~\citep{pan2025trainingsoftwareengineeringagents} and RepoNavigator~\citep{zhang2025one} further validate RL for software engineering agents, while RepoSearch-R1~\citep{li2025empoweringrepoqaagentbasedreinforcement} explores MCTS-driven self-training. Prior studies also highlight the importance of stabilizing training and controlling failure modes in group-based RL~\citep{sid2025preview}.

\subsection{Multi-agent and modular systems}
Modular architectures decompose complex software engineering tasks into specialized components. MASAI~\citep{arora2024masaimodulararchitecturesoftwareengineering} employs multiple subagents for issue reproduction, localization, fixing, and ranking. In industry practice, specialized subagents for retrieval and context gathering have also become common~\citep{cognition2025swegrep,morph2025warpgrep,cherny2025claudecode,heule2025semsearch}. Such modularity motivates training dedicated localization components that can be integrated with downstream editing agents.

% \subsection{Positioning}
% Our work targets multi-level localization (file/class/function) with a lightweight tool scaffold and trains a dedicated localization subagent directly with RL. Compared to structure-heavy approaches (e.g., graph-guided localization) and general repository-level agents, we emphasize a simple, reproducible recipe with multi-level F1 rewards and stability techniques, and release the full training and evaluation pipeline.
